LANCE is a language-model agent designed to conduct and evaluate security audits
of IoT/OT networks. The main objective of this internship was to develop an
agentic HMoE architecture in which a quantized foundation model is shared by
QLoRA adapters specialized in campaign coordination, reconnaissance,
vulnerability analysis, and exploitation. A data-processing chain converts
execution traces into validated examples, while a deterministic router selects
the expert assigned to each phase through a common API. The execution pipeline
controls transitions, tools, output schemas, and evidence. In parallel, the
experimental harness deploys reproducible scenarios and measures behavior during
campaigns. The evaluation combines deterministic metrics with constrained
semantic assessment, without delegating final score computation to a language
model. Context-management mechanisms and a weighted attack graph with an impact
cone complement the architecture. Initial results reveal a gap between produced
findings and fully evidenced attack paths. They must still be confirmed through
repeated campaigns comparing specialized experts with a general-purpose
baseline under identical conditions. This work therefore treats model
specialization, agent orchestration, and experimental measurement as inseparable
parts of a verifiable cybersecurity system.

\vspace{1em}
\noindent\textbf{Keywords:} LLM agent, IoT/OT cybersecurity, HMoE, QLoRA,
agentic pipeline, reproducible benchmark.
